\documentclass[conference]{IEEEtran}
\IEEEoverridecommandlockouts

% ============================================================
% Pacotes básicos
% ============================================================
\usepackage[utf8]{inputenc}
\usepackage[T1]{fontenc}
\usepackage[brazil]{babel}
\usepackage{lmodern}
\usepackage{microtype}
\usepackage{graphicx}
\usepackage{amsmath,amssymb}
\usepackage{booktabs}
\usepackage{tabularx}
\usepackage{array}
\usepackage{multirow}
\usepackage[table]{xcolor}
\usepackage{hyperref}
\usepackage{cite}
\usepackage{enumitem}
\usepackage{listings}
\usepackage{float}
\usepackage{stfloats}
\usepackage{tikz}
\usepackage{url}
\usetikzlibrary{arrows.meta,positioning,shapes.geometric,fit}

% ============================================================
% Configurações de estilo
% ============================================================
\hypersetup{
    colorlinks=true,
    linkcolor=black,
    citecolor=black,
    urlcolor=blue
}

\definecolor{headerblue}{RGB}{27,79,114}
\definecolor{rowblue}{RGB}{238,244,249}
\definecolor{codegray}{RGB}{245,245,245}
\definecolor{codeblue}{RGB}{0,0,160}
\definecolor{lightgray}{RGB}{248,248,248}
\definecolor{softblue}{RGB}{231,241,252}

\lstset{
    backgroundcolor=\color{codegray},
    basicstyle=\scriptsize\ttfamily,
    keywordstyle=\color{codeblue}\bfseries,
    breaklines=true,
    breakatwhitespace=false,
    columns=flexible,
    keepspaces=true,
    showstringspaces=false,
    frame=single,
    framerule=0.3pt,
    captionpos=b,
    xleftmargin=0.5em,
    xrightmargin=0.5em
}

\newcolumntype{Y}{>{\raggedright\arraybackslash}X}
\renewcommand{\arraystretch}{1.18}
\setlength{\tabcolsep}{3pt}
\setlist[itemize]{leftmargin=1.2em,itemsep=2pt,topsep=2pt,parsep=0pt}
\setlist[enumerate]{leftmargin=1.5em,itemsep=2pt,topsep=2pt,parsep=0pt}
\emergencystretch=3em
\sloppy

\newcommand{\safeincludegraphics}[2][]{%
  \IfFileExists{#2}{\includegraphics[#1]{#2}}{%
    \fbox{\parbox{0.92\linewidth}{\centering\vspace{0.35cm}
    \textbf{Imagem não encontrada}\par\small\texttt{\detokenize{#2}}\vspace{0.35cm}}}}
}

% ============================================================
\begin{document}

\title{SCC0270 -- Relatório Final do Projeto\\
\large Classificação de Insetos em Armadilhas Adesivas com\\
Redes Neurais Convolucionais}

\author{
  \IEEEauthorblockN{
    Felipe da Costa Coqueiro -- NUSP: 11781361 \\
    Guilherme Bertinati Silvestre -- NUSP: 14575683 \\
    Fernando Alee Suaiden -- NUSP: 12680836 \\
  }
  \IEEEauthorblockA{
    \textit{SCC0270 -- Redes Neurais e Aprendizado Profundo}\\
    \textit{Docente: Ricardo Cerri}\\
    \textit{Instituto de Ciências Matemáticas e de Computação -- ICMC/USP}\\
    São Carlos, Brasil, 2026
  }
}

\maketitle

% ============================================================
\begin{abstract}
Este trabalho apresenta uma solução de classificação binária de insetos
capturados em armadilhas adesivas amarelas do dataset público EMBRAPA.
O objetivo é distinguir a mosca-branca \textit{Bemisia tabaci} (classe \texttt{WF})
de outros insetos presentes nas armadilhas (classe \texttt{MR}), utilizando
recortes de bounding boxes já disponíveis como entradas de classificação.
A divisão dos dados foi realizada por imagem original para evitar vazamento
de informação entre partições. Foram treinados três modelos: um baseline
geométrico baseado em Regressão Logística com atributos de tamanho, um
Perceptron Multicamadas (MLP) e uma Rede Neural Convolucional (CNN).
A métrica principal é a Área sob a Curva Precisão-Revocação
(\textit{Average Precision}, AP). Os resultados mostram AP de \textbf{0,9994}
para o baseline geométrico, \textbf{0,9963} para o MLP e \textbf{1,0000}
para a CNN. A análise revela que atributos geométricos simples já são
altamente discriminativos neste dataset, mas a CNN apresenta vantagens
operacionais adicionais, especialmente na precisão da classe minoritária
\texttt{MR} (\textbf{0,9885} vs. \textbf{0,6457} do MLP no limiar 0,5).
\end{abstract}

\begin{IEEEkeywords}
Classificação de insetos, Redes Neurais Convolucionais, MLP,
Armadilhas adesivas, EMBRAPA, Average Precision, Aprendizado Profundo.
\end{IEEEkeywords}

% ============================================================
\section{Introdução}

O monitoramento de pragas agrícolas é uma tarefa crítica para a
produtividade e sustentabilidade da agricultura. A mosca-branca
\textit{Bemisia tabaci}, conhecida como \textit{whitefly} (WF), é uma
das pragas mais prejudiciais à produção de diversas culturas no Brasil
e no mundo, sendo responsável por danos diretos e pela transmissão de
vírus fitopatogênicos~\cite{barbedo2014}. O uso de armadilhas adesivas
amarelas é um método tradicional de monitoramento, mas a contagem e
identificação manual dos insetos capturados é trabalhosa e sujeita a
erros humanos.

Nesse contexto, técnicas de visão computacional e aprendizado profundo
surgem como alternativas promissoras para automatizar a identificação de
insetos em imagens digitalizadas de armadilhas~\cite{jige2017}. Este
projeto aborda o problema de classificação binária de recortes de insetos
em duas classes: \texttt{WF} (mosca-branca \textit{Bemisia tabaci}) e
\texttt{MR} (demais insetos presentes nas armadilhas).

O dataset utilizado foi disponibilizado pela EMBRAPA e contém imagens de
armadilhas adesivas com anotações em formato Pascal VOC. A partir dessas
anotações, foram extraídos recortes individuais de insetos, formando o
conjunto de dados utilizado neste trabalho.

As questões de pesquisa que orientam este projeto são:
\begin{itemize}
    \item \textbf{RQ1:} Qual é o poder discriminativo de atributos geométricos
    simples (largura, altura, área e razão de aspecto) neste dataset?
    \item \textbf{RQ2:} Qual é a diferença de desempenho entre MLP e CNN
    quando avaliados por Average Precision?
    \item \textbf{RQ3:} Como os modelos se comportam em relação à classe
    minoritária \texttt{MR} no limiar de decisão 0,5?
\end{itemize}

% ============================================================
\section{Trabalhos Relacionados}

A classificação automática de insetos em imagens de armadilhas tem sido
investigada com diversas abordagens. Barbedo et al.~\cite{barbedo2014}
discutem os desafios de contagem automatizada de insetos em imagens de
armadilhas cromáticas, apontando a variabilidade de escala, postura e
oclusão como principais dificuldades.

Ji et al.~\cite{jige2017} propõem o uso de CNNs para reconhecimento de
pragas agrícolas, demonstrando que características locais aprendidas
por convoluções são mais robustas à deformação de ponto de vista do que
descritores manuais. Barré et al.~\cite{bar2019} avaliam o potencial de
redes profundas para identificação de pestes em folhas, reportando ganhos
consistentes sobre métodos clássicos de extração de features.

Davis e Goadrich~\cite{davis2006} discutem as propriedades das curvas
Precisão-Revocação (PR) em relação às curvas ROC, argumentando que PR-AUC
é mais informativa em problemas com classes desbalanceadas, justificando
a adoção de Average Precision como métrica principal neste trabalho.

% ============================================================
\section{Dataset}

\subsection{Descrição Geral}

O dataset EMBRAPA de moscas brancas contém 284 imagens originais de
armadilhas adesivas amarelas anotadas em formato Pascal VOC (XML).
A partir das anotações de bounding boxes, foram extraídos \textbf{7.637
recortes}, distribuídos em duas classes:

\begin{itemize}
    \item \texttt{WF} (\textit{whitefly}): \textbf{6.350 recortes} (83,1\%)
    \item \texttt{MR} (demais insetos): \textbf{1.287 recortes} (16,9\%)
\end{itemize}

O conjunto é fortemente desbalanceado, com proporção WF:MR de
aproximadamente 4,9:1. Esse desbalanceamento foi tratado com amostragem
ponderada (\texttt{WeightedRandomSampler}) durante o treinamento.

\begin{figure}[H]
  \centering
  \safeincludegraphics[width=\columnwidth]{results/distribuicao_recortes.png}
  \caption{Distribuição das classes e dimensões dos recortes.
  WF é majoritário e tende a ter dimensões menores que MR.}
  \label{fig:distribuicao}
\end{figure}

\subsection{Análise Exploratória}

A análise estatística das dimensões revelou uma diferença sistemática
entre as classes (Tabela~\ref{tab:stats}): recortes \texttt{WF} tendem
a ser menores, enquanto recortes \texttt{MR} são maiores. Essa distinção
geométrica é o principal fator que torna o baseline geométrico tão
competitivo neste dataset.

\begin{table}[H]
\centering
\caption{Estatísticas dimensionais dos recortes por classe}
\label{tab:stats}
\footnotesize
\begin{tabularx}{\columnwidth}{lcccc}
\toprule
\textbf{Classe} & \textbf{Largura média} & \textbf{Altura média} &
\textbf{Área média} & \textbf{Razão média} \\
\midrule
\texttt{WF} & $\approx$\,27\,px & $\approx$\,26\,px & $\approx$\,735\,px$^2$ & $\approx$\,1,06 \\
\texttt{MR} & $\approx$\,72\,px & $\approx$\,63\,px & $\approx$\,5.500\,px$^2$ & $\approx$\,1,18 \\
\bottomrule
\end{tabularx}
\end{table}

\begin{figure}[H]
  \centering
  \safeincludegraphics[width=\columnwidth]{results/amostras_wf.png}
  \caption{Exemplos de recortes da classe \texttt{WF} (mosca-branca
  \textit{Bemisia tabaci}).}
  \label{fig:amostras-wf}
\end{figure}

\begin{figure}[H]
  \centering
  \safeincludegraphics[width=\columnwidth]{results/amostras_mr.png}
  \caption{Exemplos de recortes da classe \texttt{MR} (demais insetos).}
  \label{fig:amostras-mr}
\end{figure}

% ============================================================
\section{Metodologia}

\subsection{Divisão sem Vazamento de Informação}

O dataset contém recortes provenientes de 284 imagens originais.
Se a divisão fosse feita diretamente por recorte, recortes da mesma
imagem poderiam aparecer simultaneamente em treino e teste, produzindo
\textit{data leakage} e métricas de teste otimistas~\cite{barbedo2014}.

Para evitar esse problema, a divisão foi feita \textbf{por imagem original},
utilizando \texttt{GroupShuffleSplit} da biblioteca scikit-learn.
O identificador numérico presente no nome do arquivo (ex.: \texttt{1000}
em \texttt{1000\_WF\_001.jpg}) é usado como grupo. As proporções obtidas foram:

\begin{table}[H]
\centering
\caption{Partições do dataset (por imagem original)}
\label{tab:split}
\footnotesize
\begin{tabularx}{\columnwidth}{lcccc}
\toprule
\textbf{Partição} & \textbf{Recortes} & \textbf{WF} & \textbf{MR} &
\textbf{Imagens origem} \\
\midrule
Treino    & $\approx$70\% & -- & -- & 0 em comum com val/teste \\
Validação & $\approx$15\% & -- & -- & 0 em comum com treino/teste \\
Teste     & $\approx$15\% & -- & -- & 0 em comum com treino/val \\
\bottomrule
\end{tabularx}
\end{table}

A verificação explícita confirmou que nenhum grupo de origem se repete
entre as partições (\texttt{grupos\_em\_comum = 0} em todas as
comparações).

\subsection{Pré-processamento}

As imagens foram processadas de forma distinta para MLP e CNN:

\subsubsection{MLP}
Redimensionamento direto para $32 \times 32$ pixels, sem padding,
seguido de conversão para tensor. O vetor de entrada tem dimensão
$32 \times 32 \times 3 = 3.072$.

\subsubsection{CNN}
\begin{enumerate}
    \item \textbf{PadToSquare}: padding com amarelo \texttt{(255,255,0)}
    para tornar a imagem quadrada, preservando a razão de aspecto original
    e evitando distorções geométricas. O fundo amarelo mimetiza a cor
    da armadilha adesiva, minimizando artefatos de borda.
    \item \textbf{Resize} para $64 \times 64$ pixels.
    \item \textbf{Data augmentation} (apenas no treino): espelhamento
    horizontal aleatório e rotação de até $10°$.
    \item \textbf{Normalização}: média $[0{,}5,\;0{,}5,\;0{,}2]$,
    desvio-padrão $[0{,}25,\;0{,}25,\;0{,}1]$ por canal.
\end{enumerate}

\subsection{Balanceamento de Classes}

O desbalanceamento entre \texttt{WF} (6.350) e \texttt{MR} (1.287)
foi tratado com \texttt{WeightedRandomSampler}, que atribui pesos
inversamente proporcionais à frequência de cada classe. Com isso, a
rede vê, em cada época, amostras das duas classes com frequência
aproximadamente equilibrada.

% ============================================================
\section{Modelos}

\subsection{Baseline Geométrico}

O baseline geométrico é um \textit{pipeline} composto por
\texttt{StandardScaler} e \texttt{LogisticRegression}
(\texttt{max\_iter=5000}), treinado sobre quatro atributos numéricos:
largura, altura, área e razão de aspecto dos recortes. Não envolve
processamento visual. Seu propósito é quantificar o quanto de
discriminação pode ser obtido apenas com a escala do bounding box.

\subsection{MLP (Perceptron Multicamadas)}

\begin{table}[H]
\centering
\caption{Arquitetura do MLPClassifier}
\label{tab:mlp}
\footnotesize
\begin{tabularx}{\columnwidth}{lcc}
\toprule
\textbf{Camada} & \textbf{Saída} & \textbf{Detalhes} \\
\midrule
Linear + ReLU & 256 & Entrada: 3.072 \\
Dropout(0,3) & 256 & -- \\
Linear + ReLU & 64 & -- \\
Linear & 1 & Logit binário \\
\midrule
\multicolumn{3}{c}{\textbf{Parâmetros totais: 803.201}} \\
\bottomrule
\end{tabularx}
\end{table}

O MLP recebe a imagem inteiramente achatada, descartando qualquer
estrutura espacial. Serve como referência de desempenho para demonstrar
o ganho obtido pela CNN.

\subsection{CNN (Rede Neural Convolucional)}

\begin{table}[H]
\centering
\caption{Arquitetura do CNNClassifier}
\label{tab:cnn}
\footnotesize
\begin{tabularx}{\columnwidth}{lcc}
\toprule
\textbf{Bloco} & \textbf{Canais} & \textbf{Operações} \\
\midrule
Bloco 1 & 24 & Conv3×3 → BN → ReLU → Conv3×3 → BN → ReLU → MaxPool2 \\
Bloco 2 & 48 & Conv3×3 → BN → ReLU → Conv3×3 → BN → ReLU → MaxPool2 \\
Bloco 3 & 96 & Conv3×3 → BN → ReLU → MaxPool2 \\
Cabeça  & 1  & AdaptiveAvgPool(1) → Flatten → Dropout(0,3) → Linear(96→1) \\
\midrule
\multicolumn{3}{c}{\textbf{Parâmetros totais: 79.225 ($\approx$10$\times$ menos que o MLP)}} \\
\bottomrule
\end{tabularx}
\end{table}

O uso de \texttt{BatchNorm2d} em cada bloco convolucional estabiliza o
treinamento e reduz a sensibilidade à taxa de aprendizado. O
\texttt{AdaptiveAvgPool2d(1)} torna a arquitetura invariante ao tamanho
de entrada, consolidando o mapa de ativações em um vetor de 96
dimensões antes da camada de classificação.

% ============================================================
\section{Configuração Experimental}

\subsection{Treinamento}

Todos os modelos neurais foram treinados com os seguintes parâmetros:

\begin{table}[H]
\centering
\caption{Hiperparâmetros de treinamento}
\label{tab:hiper}
\footnotesize
\begin{tabularx}{\columnwidth}{lll}
\toprule
\textbf{Hiperparâmetro} & \textbf{MLP} & \textbf{CNN} \\
\midrule
Épocas          & 4 & 4 \\
Taxa de aprendizado & $10^{-3}$ & $8 \times 10^{-4}$ \\
Otimizador      & Adam & Adam \\
Função de perda & BCEWithLogitsLoss & BCEWithLogitsLoss \\
Batch size treino & 128 & 128 \\
Batch size val/teste & 256 & 256 \\
Dispositivo     & CPU/CUDA & CPU/CUDA \\
\bottomrule
\end{tabularx}
\end{table}

A melhor época é selecionada por \texttt{val\_ap} (AP de validação),
e o estado do modelo nessa época é restaurado para avaliação no teste.

\subsection{Métrica Principal}

A métrica principal adotada é a \textbf{Average Precision} (AP), definida
como a área sob a curva Precisão-Revocação:
\begin{equation}
    \mathrm{AP} = \sum_{k} (R_k - R_{k-1}) \cdot P_k
\end{equation}
onde $P_k$ e $R_k$ são a precisão e a revocação no $k$-ésimo ponto da
curva. AP é preferível à AUC-ROC em problemas desbalanceados, pois é
mais sensível ao desempenho na classe positiva minoritária~\cite{davis2006}.

\subsection{Reprodutibilidade}

Para garantir reprodutibilidade, foram fixadas sementes aleatórias em
\texttt{RANDOM\_SEED = 42} para \texttt{random}, \texttt{numpy} e
\texttt{torch}. Os resultados podem variar ligeiramente em execuções
com GPU devido ao não-determinismo de operações CUDA.

% ============================================================
\section{Resultados}

\subsection{Histórico de Treinamento}

\begin{figure}[H]
  \centering
  \safeincludegraphics[width=\columnwidth]{results/mlp_historico.png}
  \caption{Histórico de treinamento do MLP: loss por época e AP de validação.}
  \label{fig:mlp-hist}
\end{figure}

\begin{figure}[H]
  \centering
  \safeincludegraphics[width=\columnwidth]{results/cnn_historico.png}
  \caption{Histórico de treinamento da CNN: loss por época e AP de validação.}
  \label{fig:cnn-hist}
\end{figure}

A CNN apresenta convergência mais rápida e AP de validação mais estável
ao longo das épocas. O MLP mostra maior oscilação na perda de validação,
indicando menor estabilidade de treinamento mesmo com poucas épocas.

\subsection{Curvas Precisão-Revocação}

\begin{figure}[H]
  \centering
  \safeincludegraphics[width=\columnwidth]{results/baseline_geometrico_pr.png}
  \caption{Curva PR do baseline geométrico no conjunto de teste (AP = 0,9994).}
  \label{fig:pr-geom}
\end{figure}

\begin{figure}[H]
  \centering
  \safeincludegraphics[width=\columnwidth]{results/mlp_pr.png}
  \caption{Curva PR do MLP no conjunto de teste (AP = 0,9963).}
  \label{fig:pr-mlp}
\end{figure}

\begin{figure}[H]
  \centering
  \safeincludegraphics[width=\columnwidth]{results/cnn_pr.png}
  \caption{Curva PR da CNN no conjunto de teste (AP = 1,0000).}
  \label{fig:pr-cnn}
\end{figure}

As três curvas PR apresentam AP muito próxima de 1. A diferença mais
relevante está na \textit{forma} das curvas: a CNN mantém precisão próxima
a 1,0 em praticamente toda a faixa de revocação, enquanto o MLP apresenta
quedas de precisão em limiares mais baixos.

\subsection{Matrizes de Confusão}

\begin{figure}[H]
  \centering
  \safeincludegraphics[width=0.8\columnwidth]{results/baseline_geometrico_confusao.png}
  \caption{Matriz de confusão do baseline geométrico (limiar 0,5).}
  \label{fig:cm-geom}
\end{figure}

\begin{figure}[H]
  \centering
  \safeincludegraphics[width=0.8\columnwidth]{results/mlp_confusao.png}
  \caption{Matriz de confusão do MLP (limiar 0,5).}
  \label{fig:cm-mlp}
\end{figure}

\begin{figure}[H]
  \centering
  \safeincludegraphics[width=0.8\columnwidth]{results/cnn_confusao.png}
  \caption{Matriz de confusão da CNN (limiar 0,5).}
  \label{fig:cm-cnn}
\end{figure}

\subsection{Tabela Comparativa de Resultados}

\begin{table*}[t]
\centering
\caption{Comparação dos modelos no conjunto de teste}
\label{tab:resultados}
\renewcommand{\arraystretch}{1.25}
\begin{tabularx}{\textwidth}{p{2.2cm} r r r r r r}
\rowcolor{headerblue}
\color{white}\textbf{Modelo} &
\color{white}\textbf{Parâmetros} &
\color{white}\textbf{Test AP} &
\color{white}\textbf{Precisão MR@0,5} &
\color{white}\textbf{Revocação MR@0,5} &
\color{white}\textbf{Acurácia@0,5} &
\color{white}\textbf{Test Loss} \\
\rowcolor{rowblue}
CNN & 79.225 & \textbf{1,0000} & \textbf{0,9885} & \textbf{1,0000} & \textbf{0,9982} & \textbf{0,0282} \\
Baseline Geométrico & n/a & 0,9994 & 0,9429 & 0,9593 & 0,9850 & -- \\
\rowcolor{rowblue}
MLP & 803.201 & 0,9963 & 0,6457 & 0,9535 & 0,9134 & 0,1901 \\
\bottomrule
\end{tabularx}
\end{table*}

Os resultados estão sumarizados na Tabela~\ref{tab:resultados}.
Os três modelos atingem AP superior a 0,99, mas diferem
significativamente no comportamento operacional no limiar 0,5.
A CNN apresenta precisão de 0,9885 para \texttt{MR}, enquanto o MLP
obtém apenas 0,6457, indicando que o MLP classifica incorretamente
grande parte dos insetos não-WF como WF nesse limiar.

% ============================================================
\section{Discussão}

\subsection{Resposta à RQ1 -- Poder Discriminativo Geométrico}

O baseline geométrico com apenas quatro atributos de tamanho atingiu
AP = 0,9994 no conjunto de teste. Esse resultado confirma a hipótese
de que o tamanho do bounding box é um \textit{shortcut} extremamente
forte neste dataset: mosca-branca e outros insetos diferem
sistematicamente em escala, e a Regressão Logística consegue explorar
esse sinal com altíssima eficácia.

Essa observação tem implicação metodológica importante: qualquer modelo
que não supere o baseline geométrico estaria adicionando complexidade
sem ganho. Tanto o MLP quanto a CNN superam o baseline em AP, mas
especialmente em precisão para \texttt{MR}, onde o baseline apresenta
0,9429 e a CNN atinge 0,9885.

\subsection{Resposta à RQ2 -- MLP vs. CNN}

A CNN supera o MLP em todas as métricas avaliadas, com uma diferença
particularmente expressiva na precisão da classe minoritária:
\textbf{+0,3428} (0,9885 -- 0,6457). Ao mesmo tempo, a CNN possui
$\approx$10 vezes menos parâmetros (79.225 vs. 803.201), demonstrando
que localidade indutiva das convoluções é mais eficiente do que a
capacidade bruta de parâmetros de uma rede totalmente conectada para
dados com estrutura espacial.

O padding com amarelo antes do redimensionamento foi uma decisão
importante: sem ele, recortes com razões de aspecto extremas (insetos
muito alongados ou muito compactos) seriam esticados artificialmente,
distorcendo os padrões visuais que a CNN deve aprender.

\subsection{Resposta à RQ3 -- Comportamento no Limiar 0,5}

No limiar fixo 0,5, a CNN produz apenas falsos positivos mínimos
para \texttt{MR} (revocação 1,0000), enquanto o MLP deixa passar
cerca de 35\% das instâncias \texttt{MR} como \texttt{WF}
(precisão = 0,6457). Em um cenário operacional real de monitoramento
de pragas, erros no sentido de não detectar \texttt{MR}
(falsos negativos) têm custo menor do que erros na direção oposta.
A CNN apresenta melhor equilíbrio entre precisão e revocação.

\subsection{Limitações}

\begin{itemize}
    \item \textbf{Shortcut geométrico}: o dataset possui forte
    correlação entre tamanho do bounding box e classe, o que pode
    inflar métricas em relação a cenários reais mais difíceis, onde
    WF e MR poderiam ter tamanhos similares.
    \item \textbf{Épocas limitadas}: com 4 épocas, os modelos podem
    não ter convergido totalmente. Experimentos com mais épocas e
    \textit{early stopping} poderiam revelar desempenho distinto.
    \item \textbf{Generalização}: o dataset provém de uma única fonte
    (EMBRAPA), o que pode não representar a diversidade de condições
    de captura em campo.
    \item \textbf{Transferência de aprendizado}: modelos pré-treinados
    como ResNet18 ou MobileNetV3 não foram avaliados neste trabalho,
    mas poderiam oferecer melhor generalização com poucos dados.
\end{itemize}

% ============================================================
\section{Conclusão}

Este trabalho apresentou uma solução de classificação binária de insetos
em armadilhas adesivas utilizando o dataset EMBRAPA. A divisão por imagem
original via \texttt{GroupShuffleSplit} evitou vazamento de informação,
garantindo que a avaliação no conjunto de teste refletisse corretamente
a capacidade de generalização dos modelos.

Os três modelos avaliados atingiram AP superior a 0,99, indicando que
o dataset é altamente separável. O baseline geométrico demonstrou que
atributos de tamanho já contêm informação discriminativa muito forte.
O MLP, embora com AP elevada, apresentou comportamento operacional
inferior no limiar 0,5, especialmente para a classe minoritária
\texttt{MR}. A CNN obteve o melhor resultado geral:
AP = \textbf{1,0000}, precisão MR = \textbf{0,9885},
revocação MR = \textbf{1,0000} e acurácia = \textbf{0,9982},
com apenas 79.225 parâmetros.

A metodologia adotada -- divisão sem vazamento, tratamento de
desbalanceamento, padding com preservação de proporção e
\texttt{BatchNorm} -- foi fundamental para garantir a qualidade dos
resultados. Como trabalhos futuros, sugere-se avaliar transferência
de aprendizado com arquiteturas pré-treinadas, aumentar o número de
épocas com \textit{early stopping} baseado em AP de validação e
explorar datasets com maior variabilidade de condições de captura.

% ============================================================
\begin{thebibliography}{00}

\bibitem{barbedo2014}
J. G. A. Barbedo,
``Using digital image processing for counting whiteflies on soybean leaves,''
\textit{Journal of Asia-Pacific Entomology}, vol. 17, no. 4,
pp. 685--694, 2014.

\bibitem{jige2017}
G. Ji, L. Zhu, X. Gu, and C. Lu,
``Automatic pest detection and recognition in field crops using deep
convolutional neural network,''
in \textit{Proc. Int. Conf. Image and Graphics (ICIG)},
Springer, 2017, pp. 519--530.

\bibitem{bar2019}
P. Barré, B. C. Stöver, K. F. Müller, and V. Steinhage,
``LeafNet: A computer vision system for automatic plant species
identification,''
\textit{Ecological Informatics}, vol. 40, pp. 50--56, 2017.

\bibitem{davis2006}
J. Davis and M. Goadrich,
``The relationship between Precision-Recall and ROC curves,''
in \textit{Proc. 23rd Int. Conf. Machine Learning (ICML)},
ACM, 2006, pp. 233--240.

\end{thebibliography}

\end{document}